% GENERATED FILE. Edit canonical lessons, then run npm run generate:books.
% canonical-source-hash: fnv1a64:fa1a7b4a83381919
% canonical-lessons: AR-C409-qanun, AR-C409-tashri, AR-C409-band, AR-C409-nassa, AR-C409-sharraa

\chapter{Laws}
\label{ch:ar-laws}
\hlchaptermodality{409}

\begin{chapteropening}
\textbf{By the end of this chapter:} \emph{I can talk about a law, legislation and a clause, and say what a law stipulates or who legislates.}

Complete the last lesson of chapter 409: I can talk about a law, legislation and a clause, and say what a law stipulates or who legislates.
\end{chapteropening}

\section[\texorpdfstring{\ar{قانون} (qānūn)}{قانون (qānūn)}]{\ar{قانون} (qānūn) — a law}

\label{lesson:AR-C409-qanun}



\begin{quote}
Before the new one: say the Arabic for well-being, then the Arabic for care.
\end{quote}

\subsection*{You'll want to know: \ar{قانون}}
\textbf{\ar{قانون}} — \emph{qānūn} — \textquotedblleft{}a law\textquotedblright{}.

\begin{quote}
\ar{يمنع} \ar{القانون} \ar{التدخين} \ar{في} \ar{المطاعم}.
\end{quote}

\begin{quote}
\emph{yamnaʿu al-qānūn at-tadkhīn fī al-maṭāʿim.}
\end{quote}

\begin{quote}
The law bans smoking in restaurants.
\end{quote}

Plural: \ar{قوانين}, \emph{qawānīn}.

\begin{cousinweb}
\textbf{\ar{قانون}} (\emph{qānūn}) was borrowed from Greek \emph{kanōn}, \textquotedblleft{}a measuring rod, a rule\textquotedblright{} (EtymDB). Arabic gave it the sense of a law.
\end{cousinweb}

\begin{grammarlens}[title={What you've built}]
One more word for saying what happens if.
\end{grammarlens}

\subsection*{Your turn}
\begin{itemize}
\raggedright
  \item \emph{Say it:} \emph{qānūn}
  \item \emph{Say it:} \emph{qānūn}, once more
  \item \emph{Say it:} \emph{yamnaʿu al-qānūn at-tadkhīn fī al-maṭāʿim.} — The law bans smoking in restaurants.
  \item \emph{Recall:} say the Arabic for well-being, then the Arabic for care, then say \emph{qānūn} again
\end{itemize}

\begin{tcolorbox}[breakable,colback=teal!4,colframe=teal!35!black,title={Before you move on}]
What does \ar{قانون} mean? (\textquotedblleft{}A law\textquotedblright{}.) Say it once more.
\end{tcolorbox}
\section[\texorpdfstring{\ar{تشريع} (tashrīʿ)}{تشريع (tashrīʿ)}]{\ar{تشريع} (tashrīʿ) — legislation}

\label{lesson:AR-C409-tashri}



\begin{quote}
Before the new one: say the Arabic for care, then the Arabic for a law.
\end{quote}

\subsection*{You'll want to know: \ar{تشريع}}
\textbf{\ar{تشريع}} — \emph{tashrīʿ} — \textquotedblleft{}legislation\textquotedblright{}.

\begin{quote}
\ar{تحتاج} \ar{البلاد} \ar{إلى} \ar{تشريع} \ar{جديد} \ar{للعمل}.
\end{quote}

\begin{quote}
\emph{taḥtāju al-bilād ilā tashrīʿ jadīd li-l-ʿamal.}
\end{quote}

\begin{quote}
The country needs new labour legislation.
\end{quote}

\begin{cousinweb}
\textbf{\ar{تشريع}} (\emph{tashrīʿ}) is built on the native Arabic root \emph{sh-r-ʿ}, \textquotedblleft{}to lay down a way\textquotedblright{} (Lane). \emph{Tashrīʿ} is the Form II verbal noun, laying down a way for others to follow: legislation.
\end{cousinweb}

\begin{grammarlens}[title={What you've built}]
One more word for saying what happens if.
\end{grammarlens}

\subsection*{Your turn}
\begin{itemize}
\raggedright
  \item \emph{Say it:} \emph{tashrīʿ}
  \item \emph{Say it:} \emph{tashrīʿ}, once more
  \item \emph{Say it:} \emph{taḥtāju al-bilād ilā tashrīʿ jadīd li-l-ʿamal.} — The country needs new labour legislation.
  \item \emph{Recall:} say the Arabic for care, then the Arabic for a law, then say \emph{tashrīʿ} again
\end{itemize}

\begin{tcolorbox}[breakable,colback=teal!4,colframe=teal!35!black,title={Before you move on}]
What does \ar{تشريع} mean? (\textquotedblleft{}Legislation\textquotedblright{}.) Say it once more.
\end{tcolorbox}
\section[\texorpdfstring{\ar{بند} (band)}{بند (band)}]{\ar{بند} (band) — a clause}

\label{lesson:AR-C409-band}



\begin{quote}
Before the new one: say the Arabic for a law, then the Arabic for legislation.
\end{quote}

\subsection*{You'll want to know: \ar{بند}}
\textbf{\ar{بند}} — \emph{band} — \textquotedblleft{}a clause\textquotedblright{}.

\begin{quote}
\ar{قرأ} \ar{المحامي} \ar{كل} \ar{بند} \ar{في} \ar{العقد}.
\end{quote}

\begin{quote}
\emph{qaraʾa al-muḥāmī kulla band fī al-ʿaqd.}
\end{quote}

\begin{quote}
The lawyer read every clause in the contract.
\end{quote}

Plural: \ar{بنود}, \emph{bunūd}.

\begin{cousinweb}
\textbf{\ar{بند}} (\emph{band}) was borrowed from Persian \emph{band}, \textquotedblleft{}a tie, a bond\textquotedblright{} (EtymDB). In Arabic it is an article or clause of a text.
\end{cousinweb}

\begin{grammarlens}[title={What you've built}]
One more word for saying what happens if.
\end{grammarlens}

\subsection*{Your turn}
\begin{itemize}
\raggedright
  \item \emph{Say it:} \emph{band}
  \item \emph{Say it:} \emph{band}, once more
  \item \emph{Say it:} \emph{qaraʾa al-muḥāmī kulla band fī al-ʿaqd.} — The lawyer read every clause in the contract.
  \item \emph{Recall:} say the Arabic for a law, then the Arabic for legislation, then say \emph{band} again
\end{itemize}

\begin{tcolorbox}[breakable,colback=teal!4,colframe=teal!35!black,title={Before you move on}]
What does \ar{بند} mean? (\textquotedblleft{}A clause\textquotedblright{}.) Say it once more.
\end{tcolorbox}
\section[\texorpdfstring{\ar{نص} (naṣṣa)}{نص (naṣṣa)}]{\ar{نص} (naṣṣa) — to stipulate in writing}

\label{lesson:AR-C409-nassa}



\begin{quote}
Before the new one: say the Arabic for legislation, then the Arabic for a clause.
\end{quote}

\subsection*{You'll want to know: \ar{نص}}
\textbf{\ar{نص}} — \emph{naṣṣa} — \textquotedblleft{}to stipulate in writing\textquotedblright{}.

\begin{quote}
\ar{ينص} \ar{القانون} \ar{على} \ar{حق} \ar{التعليم}.
\end{quote}

\begin{quote}
\emph{yanuṣṣu al-qānūn ʿalā ḥaqq at-taʿlīm.}
\end{quote}

\begin{quote}
The law stipulates the right to education.
\end{quote}

Present: \ar{ينص}, \emph{yanuṣṣ}.

It takes \ar{على}: \ar{ينص} \ar{على} …. The noun \ar{نص}, \emph{naṣṣ}, means a text.

\begin{cousinweb}
\textbf{\ar{نص}} (\emph{naṣṣa}) is built on the native Arabic root \emph{n-ṣ-ṣ}, \textquotedblleft{}to raise, to set up\textquotedblright{} (Lane). \emph{Naṣṣa} is a Form I verb; from setting a thing up plainly comes stating it in writing: to stipulate.
\end{cousinweb}

\begin{grammarlens}[title={What you've built}]
One more word for saying what happens if.
\end{grammarlens}

\subsection*{Your turn}
\begin{itemize}
\raggedright
  \item \emph{Say it:} \emph{naṣṣa}
  \item \emph{Say it:} \emph{naṣṣa}, once more
  \item \emph{Say it:} \emph{yanuṣṣu al-qānūn ʿalā ḥaqq at-taʿlīm.} — The law stipulates the right to education.
  \item \emph{Recall:} say the Arabic for legislation, then the Arabic for a clause, then say \emph{naṣṣa} again
\end{itemize}

\begin{tcolorbox}[breakable,colback=teal!4,colframe=teal!35!black,title={Before you move on}]
What does \ar{نص} mean? (\textquotedblleft{}To stipulate in writing\textquotedblright{}.) Say it once more.
\end{tcolorbox}
\section[\texorpdfstring{\ar{شرع} (sharraʿa)}{شرع (sharraʿa)}]{\ar{شرع} (sharraʿa) — to legislate}

\label{lesson:AR-C409-sharraa}



\begin{quote}
Before the new one: say the Arabic for a clause, then the Arabic for to stipulate in writing.
\end{quote}

\subsection*{You'll want to know: \ar{شرع}}
\textbf{\ar{شرع}} — \emph{sharraʿa} — \textquotedblleft{}to legislate\textquotedblright{}.

\begin{quote}
\ar{يشرع} \ar{البرلمان} \ar{القوانين}.
\end{quote}

\begin{quote}
\emph{yusharriʿu al-barlamān al-qawānīn.}
\end{quote}

\begin{quote}
Parliament makes the laws.
\end{quote}

Present: \ar{يشرع}, \emph{yusharriʿ}.

\begin{cousinweb}
\textbf{\ar{شرع}} (\emph{sharraʿa}) is built on the native Arabic root \emph{sh-r-ʿ}, \textquotedblleft{}to lay down a way\textquotedblright{} (Lane). It is Form II, to lay down a way for others: to legislate.
\end{cousinweb}

\begin{grammarlens}[title={What you've built}]
One more word for saying what happens if.
\end{grammarlens}

\subsection*{Your turn}
\begin{itemize}
\raggedright
  \item \emph{Say it:} \emph{sharraʿa}
  \item \emph{Say it:} \emph{sharraʿa}, once more
  \item \emph{Say it:} \emph{yusharriʿu al-barlamān al-qawānīn.} — Parliament makes the laws.
  \item \emph{Recall:} say the Arabic for a clause, then the Arabic for to stipulate in writing, then say \emph{sharraʿa} again
\end{itemize}

\begin{tcolorbox}[breakable,colback=teal!4,colframe=teal!35!black,title={Before you move on}]
What does \ar{شرع} mean? (\textquotedblleft{}To legislate\textquotedblright{}.) Say it once more.
\end{tcolorbox}
